% ===================== METODOLOGIA =====================
\section{Metodologia}
\label{sec:metodologia}

\subsection{Plan de Trabajo}

[Describir la metodologia del estudio como un proceso secuencial de etapas.
Cada etapa debe ser coherente con los objetivos especificos y los entregables.
Incluir tabla resumen del plan de trabajo.]

\begin{table}[htbp]
    \centering
    \caption{[Leyenda descriptiva del plan de trabajo.]}
    \label{tab:plan_trabajo}
    \setcounter{itemcount}{0}
    \begin{tabularx}{\textwidth}{@{}NlX@{}}
        \toprule
        \multicolumn{1}{c}{\textbf{Etapa}} & \textbf{Actividad} & \textbf{Descripcion} \\
        \midrule
        & [Etapa 1] & [Descripcion de la actividad.] \\
        & [Etapa 2] & [Descripcion de la actividad.] \\
        & [Etapa 3] & [Descripcion de la actividad.] \\
        & [Etapa 4] & [Descripcion de la actividad.] \\
        \bottomrule
    \end{tabularx}
\end{table}

% ─────────────────────────────────────────────────────────────────────────────
\subsection{Herramientas de Analisis}

[Describir las herramientas de software, metodos numericos o correlaciones
empleadas en el estudio. Incluir version de software, modelo matematico,
y referencia bibliografica del metodo.]

% ─────────────────────────────────────────────────────────────────────────────
\subsection{Matriz de Escenarios / Casos de Estudio}

[Si el estudio evalua multiples escenarios operacionales, presentar la matriz
completa con las combinaciones de variables y su descripcion.]

\begin{table}[htbp]
    \centering
    \caption{[Leyenda descriptiva de la matriz de escenarios.]}
    \label{tab:escenarios}
    \begin{tabularx}{\textwidth}{@{}cccccX@{}}
        \toprule
        \textbf{Escenario} & \textbf{[Var 1]} & \textbf{[Var 2]} & \textbf{[Var 3]} & \textbf{[Var 4]} & \textbf{Descripcion} \\
        \midrule
        1  & ---         & ---         & ---         & ---         & [Descripcion] \\
        2  & \checkmark  & ---         & ---         & ---         & [Descripcion] \\
        \bottomrule
    \end{tabularx}
\end{table}

\noindent\small\textit{Nota:} \checkmark~=~activo; ---~=~inactivo. [Definir abreviaturas utilizadas.]
